%% ============================================================
%%  PORTADA — Práctica 4 Alternativa
%% ============================================================
\begin{titlepage}
\centering

% ---- Banda superior institucional ----
\vspace*{1.2cm}
{\Large \textbf{Universidad de Granada}}\\[0.25cm]
{\large E.T.S. de Ingenierías Informática y de Telecomunicación}\\[0.15cm]
{\large Doble Grado en Ingeniería Informática y\\ Administración y Dirección de Empresas}\\[0.35cm]
{\normalsize Metaheurísticas \,\textbullet\, Curso 2025/2026 \,\textbullet\, 4.\textsuperscript{o} curso}

\vspace{1.6cm}

% ---- Bloque de título enmarcado ----
{\color{lcaRed}\rule{\linewidth}{2pt}}\\[0.7cm]

{\Huge \textbf{\textcolor{lcaRed}{Práctica 4}}}\\[0.25cm]
{\large\textsc{Alternativa al Examen}}\\[0.8cm]

{\LARGE \textbf{Liver Cancer Algorithm (LCA)}}\\[0.45cm]
{\large Descripción, experimentación, hibridación y extensiones}\\[0.15cm]
{\normalsize sobre el benchmark \textsc{CEC2017}}\\[0.7cm]

{\color{lcaRed}\rule{\linewidth}{2pt}}

\vspace{1.4cm}

% ---- Datos del trabajo ----
\begin{tabular}{@{}r@{\hspace{1em}}l@{}}
  \textbf{Autor}       & Ismael Sallami Moreno \\[0.35cm]
  \textbf{Metaheurística} & Liver Cancer Algorithm (Houssein \emph{et al.}, 2023) \\
  \textbf{Referencia}  & \emph{Computers in Biology and Medicine}, vol.\ 165, 107389 \\
  \textbf{DOI}         & \texttt{10.1016/j.compbiomed.2023.107389} \\[0.35cm]
  \textbf{Benchmark}   & CEC2017 — 30 funciones, $D \in \{10, 30, 50\}$, 31 ejecuciones \\
  \textbf{Repositorio} & \url{https://github.com/dmolina/cec2017real} \\[0.35cm]
  \textbf{Fecha}       & Junio de 2026 \\
\end{tabular}

\vfill
\newpage
% ---- Resumen ----
\begin{abstract}
\noindent
Este informe recoge el trabajo completo de la Práctica 4 alternativa de Metaheurísticas.
Se ha seleccionado el \textbf{Liver Cancer Algorithm (LCA)}, una metaheurística bioinspirada
propuesta en 2023 que modela el crecimiento y la propagación del carcinoma hepatocelular.
El documento se organiza en cuatro partes: (\textit{i}) descripción detallada del algoritmo y
análisis del equilibrio exploración–explotación; (\textit{ii}) estudio experimental sobre el
benchmark CEC2017 ($D\in\{10,30,50\}$, 31 ejecuciones) con análisis estadístico no paramétrico
(ranking de Friedman y test de Wilcoxon) siguiendo la metodología de \texttt{tacolab.org} y
comparando frente a DE y PSO de referencia; (\textit{iii}) diseño e implementación de un
algoritmo memético híbrido LCA\,+\,Solis-Wets; y (\textit{iv}) propuesta y evaluación de mejoras
orientadas al control de la diversidad y a la búsqueda de múltiples soluciones.
\end{abstract}

\end{titlepage}
\newpage
